\documentclass[11pt]{article}
\pagestyle{plain}

% \pgfplotsset{compat=1.18}
\usepackage[margin=1in]{geometry}
\usepackage{amsmath,amsfonts,amssymb,amsthm}
\usepackage[ruled,vlined,linesnumbered,commentsnumbered]{algorithm2e}
\SetKwInput{KwIn}{Input}
\SetKwInput{KwOut}{Output}
\SetKwInput{KwGlobal}{Global variable}
\SetKwInput{KwOracle}{Oracle}
\SetKwComment{Comment}{// }{}
\usepackage{graphicx}
\usepackage{enumerate}
\usepackage{bbm}
\usepackage{verbatim}
\usepackage{hyperref,color}
\usepackage[capitalize,nameinlink]{cleveref}
\usepackage[dvipsnames]{xcolor}
\hypersetup{
	colorlinks=true,
	pdfpagemode=UseNone,
	citecolor=OliveGreen,
	linkcolor=NavyBlue,
	urlcolor=Magenta,
	pdfstartview=FitW
}
\usepackage{appendix}
\usepackage{makecell}
\usepackage{tablefootnote}
\crefname{appsec}{Appendix}{Appendices}
\usepackage{tikz}
\usepackage{pgfplots}
\usepackage{xifthen}
\usepackage{subcaption}
\usepackage{hhline}
\usepackage{listings}

\theoremstyle{plain}
\newtheorem{theorem}{Theorem}[section]
\newtheorem{proposition}[theorem]{Proposition}
\newtheorem{lemma}[theorem]{Lemma}
\newtheorem{corollary}[theorem]{Corollary}
\newtheorem{conjecture}[theorem]{Conjecture}
\newtheorem{problem}[theorem]{Problem}
\newtheorem{claim}[theorem]{Claim}
\newtheorem{fact}[theorem]{Fact}

\theoremstyle{definition}
\newtheorem{definition}[theorem]{Definition}
\newtheorem{example}[theorem]{Example}
\newtheorem{observation}[theorem]{Observation}
\newtheorem{assumption}[theorem]{Assumption}
\newtheorem*{assumption*}{Assumption}
\newtheorem{condition}[theorem]{Condition}

\theoremstyle{remark}
\newtheorem{remark}[theorem]{Remark}

\crefname{lemma}{Lemma}{Lemmas}
\crefname{theorem}{Theorem}{Theorems}
\crefname{definition}{Definition}{Definitions}
\crefname{fact}{Fact}{Facts}
\crefname{claim}{Claim}{Claims}
\crefname{proposition}{Proposition}{Propositions}


\newcommand{\dif}{\,\mathrm{d}}
%\newcommand{\E}{\mathbb{E}}
%\newcommand{\Var}{\mathrm{Var}}
\newcommand{\Cov}{\mathrm{Cov}}
\newcommand{\MI}{\mathrm{I}}
%\newcommand{\Ent}{\mathrm{Ent}}
\newcommand{\ind}{\mathbbm{1}}
\newcommand{\allone}{\mathbf{1}}
\newcommand{\tr}{\mathrm{tr}}
\newcommand{\wt}{\mathrm{weight}}
\DeclareMathOperator*{\argmax}{arg\,max}
\newcommand{\norm}[1]{\left\lVert #1 \right\rVert}
\newcommand{\TV}[2]{\left\|#1 - #2\right\|_{\mathrm{TV}}}
\newcommand{\ceil}[1]{\left\lceil #1 \right\rceil}
\newcommand{\floor}[1]{\left\lfloor #1 \right\rfloor}
\newcommand{\ip}[2]{\left\langle #1 , #2 \right\rangle}
\newcommand{\grad}{\nabla}
\newcommand{\hessian}{\nabla^2}
\newcommand{\trans}{\intercal}
\newcommand{\diag}{\mathrm{diag}}
\newcommand{\poly}{\mathrm{poly}}
\newcommand{\dist}{\mathrm{dist}}
\newcommand{\sgn}{\mathrm{sgn}}
\newcommand{\fpras}{\mathsf{FPRAS}}
\newcommand{\fpaus}{\mathsf{FPAUS}}
\newcommand{\fptas}{\mathsf{FPTAS}}

\newcommand{\eps}{\varepsilon}

\newcommand{\N}{\mathbb{N}}
\newcommand{\Z}{\mathbb{Z}}
\newcommand{\Q}{\mathbb{Q}}
\newcommand{\R}{\mathbb{R}}
\newcommand{\F}{\mathbb{F}}
\newcommand{\C}{\mathbb{C}}

\renewcommand{\AA}{\mathcal{A}}
\newcommand{\BB}{\mathcal{B}}
\newcommand{\CC}{\mathcal{C}}
\newcommand{\DD}{\mathcal{D}}
\newcommand{\EE}{\mathcal{E}}
\newcommand{\FF}{\mathcal{F}}
\newcommand{\GG}{\mathcal{G}}
\newcommand{\HH}{\mathcal{H}}
\newcommand{\II}{\mathcal{I}}
\newcommand{\JJ}{\mathcal{J}}
\newcommand{\KK}{\mathcal{K}}
\newcommand{\LL}{\mathcal{L}}
\newcommand{\MM}{\mathcal{M}}
\newcommand{\NN}{\mathcal{N}}
\newcommand{\OO}{\mathcal{O}}
\newcommand{\PP}{\mathcal{P}}
\newcommand{\QQ}{\mathcal{Q}}
\newcommand{\RR}{\mathcal{R}}
\renewcommand{\SS}{\mathcal{S}}
\newcommand{\TT}{\mathcal{T}}
\newcommand{\UU}{\mathcal{U}}
\newcommand{\VV}{\mathcal{V}}
\newcommand{\WW}{\mathcal{W}}
\newcommand{\XX}{\mathcal{X}}
\newcommand{\YY}{\mathcal{Y}}
\newcommand{\ZZ}{\mathcal{Z}}

\newcommand{\KL}[2]{D_{\mathrm{KL}} \left( #1 \,\Vert\, #2 \right)}
\newcommand{\KLbig}[2]{D_{\mathrm{KL}} \left( #1 \,\big\Vert\, #2 \right)}
\newcommand{\KLBig}[2]{D_{\mathrm{KL}} \left( #1 \,\Big\Vert\, #2 \right)}

\newcommand{\SKL}[2]{D_{\mathrm{SKL}} \left( #1 \,\Vert\, #2 \right)}
\newcommand{\SKLbig}[2]{D_{\mathrm{SKL}} \left( #1 \,\big\Vert\, #2 \right)}
\newcommand{\SKLBig}[2]{D_{\mathrm{SKL}} \left( #1 \,\Big\Vert\, #2 \right)}

\newcommand{\qandq}{\quad\text{and}\quad}
\newcommand{\supp}{\mathrm{supp}}

\newcommand{\DKL}[2]{\-D_{\-{KL}}\left(#1 \parallel #2\right)}
\newcommand{\DTV}[2]{\-D_{\mathrm{TV}}\left({#1},{#2}\right)}
% \newcommand{\dist}{\mathrm{dist}}
\newcommand{\e}{\mathrm{e}}
\renewcommand{\epsilon}{\varepsilon}
\renewcommand{\emptyset}{\varnothing}
% \newcommand{\norm}[1]{\left\Vert#1\right\Vert}
\newcommand{\set}[1]{\left\{#1\right\}}
\newcommand{\tuple}[1]{\left(#1\right)} 
% \newcommand{\eps}{\varepsilon}
\newcommand{\inner}[2]{\left\langle #1,#2\right\rangle}
\newcommand{\tp}{\tuple}
\newcommand{\ol}{\overline}
\newcommand{\abs}[1]{\left\vert#1\right\vert}
\newcommand{\ctp}[1]{\left\lceil#1\right\rceil}
\newcommand{\ftp}[1]{\left\lfloor#1\right\rfloor}
\newcommand{\Ind}[1]{\-{Ind}_{#1}}
\newcommand{\todo}[1]{{\color{red} TODO: #1}}

\newcommand{\Cor}{\Psi^{\mathrm{sym}}}
\newcommand{\cor}{\Psi^{\-{cor}}}
\newcommand{\Inf}{\Psi}
% \newcommand{\diag}{\-{diag}}
\newcommand{\Tr}{\-{Tr}}
\newcommand{\rank}{\-{rank}}
\newcommand{\Span}{\-{span}}

\def\*#1{\boldsymbol{#1}} % Use \*A for \mathbf{A}
\def\+#1{\mathcal{#1}} % Use \+A for \mathcal{A}
\def\-#1{\mathrm{#1}} % Use \-A for \mathrm{A}
\def\=#1{\mathbb{#1}} % Use \^A for \mathbb{A}
\def\!#1{\mathfrak{#1}} % Use \!A for \mathfrak{A}

\newcommand*\diff{\mathop{}\!\mathrm{d}}

% probability
% \DeclareMathOperator*{\oPr}{\mathbf{Pr}}
\def\oPr{\mathop{\mathrm{Pr}}}
\renewcommand{\Pr}[2][]{ \ifthenelse{\isempty{#1}}
  {\oPr\left[#2\right]}
  {\oPr_{#1}\left[#2\right]} } % Use \Pr[a]{b} for \mathbf{Pr}_a[b], \Pr{b} for  \mathbf{Pr}[b]

% expectation: relies on the package "dsfont"
% \DeclareMathOperator*{\oE}{\mathbb{E}}
\def\oE{\mathop{\mathbb{E}}}
\newcommand{\E}[2][]{ \ifthenelse{\isempty{#1}}
  {\oE\left[#2\right]}
  {\oE_{#1}\left[#2\right]} }

% variance
%\DeclareMathOperator*{\oVar}{\mathbf{Var}}
\def\oVar{\mathrm{Var}}
\newcommand{\Var}[2][]{ \ifthenelse{\isempty{#1}}
  {\oVar\left[#2\right]}
  {\oVar_{#1}\left[#2\right]} }

% entropy
% \DeclareMathOperator*{\oEnt}{\mathbf{Ent}}
\def\oEnt{\mathrm{Ent}}
\newcommand{\Ent}[2][]{ \ifthenelse{\isempty{#1}}
  {\oEnt\left[#2\right]}
  {\oEnt_{#1}\left[#2\right]} }

\newcommand{\PhiEnt}[2][]{ \ifthenelse{\isempty{#1}}
  {\oEnt^\phi\left[#2\right]}
  {\oEnt^\phi_{#1}\left[#2\right]} }


\newenvironment{Exampleparagraph}[1][]
{
    \par 
    \noindent 
    \textbf{\color{blue} #1}
    \par 
    \vspace{0.5em} 
    \itshape 
}
{
    \par 
    \vspace{1em} 
}

%\input{local-macro}

\newcommand{\RED}[1]{\textcolor{red}{#1}}

\newcommand{\bern}[1]{\mathrm{Bern}\left(#1\right)}
\newcommand{\pll}{\parallel}



\title{MATH 6112. FALL 2026. ASSIGNMENT 4.}
\date{Due on November 23, 2026}
\author{{\textbf{Zejia Chen}} \\ \textbf{GTID: 904178159}}

\pgfplotsset{compat=1.18} 
\begin{document}
\maketitle
For Hermitian matrix $A \in \C^{n\times n}$, we use $\lambda_1(A)\leq \cdots \leq \lambda_n(A)$ to denote the eigenvalues of $A$.
\section*{Problem 1}
\textbf{(i)}
We have that 
\begin{equation}\label{eq:P1_1}
    \mu = \min_{x\neq 0} \frac{x^*Ax}{\norm{B^{1/2}x}_2^2} = \min_{\norm{y}_2=1} y^*B^{-\frac 12}AB^{-\frac 12}y = \lambda_1\tp{B^{-\frac 12}AB^{-\frac 12}} = \lambda_1(B^{-1}A),
\end{equation}
where the last equality follows from $\det\tp{\lambda I - B^{-1/2}AB^{-1/2}} = \det\tp{\lambda I - B^{-1}A}$.
Since $A$ is Hermitian, $B^{-1/2}AB^{-1/2}$ can be diagonalized: $B^{-1/2}AB^{-1/2} = V\Lambda V^*$, where $\Lambda$ is the diagonal matrix of eigenvalues of $B^{-1/2}AB^{-1/2}$.
Hence, $B^{-1}A = B^{-1/2}V\Lambda V^*B^{1/2}$, that is, $B^{-1}A$ is diagonalizable.
Therefore, all the eigenvalues of $B^{-1}A$ has its cooresponding eigenvector.
Suppose $B^{-1}A v = \lambda_1(B^{-1}A)v$, then $Av = \mu Bv$ and $x^*Ax = \mu x^*Bx = \mu \norm{B^{1/2}x}_2$, which gives the minimum value in \Cref{eq:P1_1}.
\\\textbf{(ii)}
Let $w_1\ldots w_n$ be the eigenvectors of $B^{-1/2}AB^{-1/2}$ with corresponding eigenvalues $\mu_1=\mu\leq\mu_2\leq\ldots \leq\mu_n$.
Suppose $w_1,\ldots w_n$ are orthonormal since $B^{-1/2}AB^{-1/2}$ is Hermitian.
Since $B^{-1}Av = \mu v$, $B^{-1/2}AB^{-1/2}(B^{1/2}v) = \mu(B^{1/2}v)$ and thus, $w_1 = B^{1/2}v$.
We have that 
\begin{align}\label{eq:P1_2}
    \nu = \min_{\substack{x\neq 0\\x\perp Bv}} \frac{x^*Ax}{\norm{B^{1/2}x}_2^2} &= \min_{\substack{\norm{y}_2=1\\y\perp B^{1/2}v}}y^*B^{-\frac 12}AB^{-\frac 12}y  = \min_{\substack{\sum_{i>1} a_i^2 = 1\\y=\sum_{i>1}a_iw_i}} y^*B^{-\frac 12}AB^{-\frac 12}y
    \\&= \min_{\sum_{i>1} a_i^2 = 1}\sum_{i=2}^n \mu_i a_i^2 = \mu_2.
\end{align}
The equality can be achieved by $B^{-1/2}w_2$ and $B^{-1/2}AB^{-1/2}w_2 = \nu w_2$ implies $A(B^{-1/2}w_2) = \nu B(B^{-1/2}w_2)$.
That is, $B^{-1/2}w_2$ is the solution of $Aw = \nu Bw$.
\section*{Problem 2}
\textbf{(a)}
Suppose $\lambda$ is an eigenvalue of $A$ and $v$ is the corresponding eigenvector with $\norm{v}_2 = 1$.
Therefore, $Av = \lambda v$ and $v^*A^* = \overline{\lambda}v^*$.
Then we have that
$$
    0 > v^*(BA+A^*B)v = 2\Re\lambda (v^*Bv) \implies \Re\lambda < 0,
$$
since $B$ is positive definite and $v^*Bv > 0$.
\\\textbf{(b)} 
Since for real $t > 0$,
$$
    \tp{e^{At}}^* = \sum_{i=0}^\infty \frac{((At)^i)^*}{i!} = \sum_{i=0}^\infty \frac{((A^*t)^i)}{i!} = e^{A^* t}, 
$$
$e^{A^*t}e^{At}$ is Hermitian, positive semidefinite and so is $B$.
For any $v\neq 0$, $e^{0\cdot A}v = v \neq 0$.
Since $t\mapsto e^{tA} v$ is continuous, there exists $\delta > 0$ such that $e^{tA}v \neq 0$ for any $t\geq \delta$.
Therefore,
$$
    \int_0^\infty v^*e^{tA^*}e^{tA}v \dif t = 
    \int_0^\infty \norm{e^{tA}v}_2^2 \dif t
    \geq \int_0^\delta \norm{e^{tA}v}_2^2 \dif t > 0,
$$
implying that $B$ is positive definite.

Observe that
$$
    \frac{\dif}{\dif t}e^{tA} = \sum_{i=0}^\infty\frac{\dif}{\dif t}\frac{t^iA^i}{i!}
    = A \sum_{i=1}^\infty \frac{(tA)^{i-1}}{(i-1)!} = A e^{tA} = e^{tA}A.
$$
Therefore,
$$
    \frac{\dif}{\dif t}\tp{e^{tA^*}e^{tA}} =  \tp{\frac{\dif}{\dif t}e^{tA^*}}e^{tA} + e^{tA^*}\tp{\frac{\dif}{\dif t} e^{tA}} = e^{tA^*}(A^*+A)e^{tA}.
$$
Applying the above equation, we conclude that
$$
    BA+A^*B = \int_0^\infty e^{tA^*} (A^*+A) e^{tA}\dif t = \left. e^{tA^*}e^{tA} \right|_{0}^{\infty} = -I \prec 0,
$$
since $e^{tA}\to 0$ as $t\to \infty$.
\\\textbf{(c)} 
We first assume that $A$ is diagonalizable.
Let $v_1,\ldots,v_n$ be linear independent eigenvectors of $A$ with eigenvalues $a_i + b_i\-i$ ($a_i < 0$ for $i\in[n]$).
Therefore, for any $i\in[n]$, by Euler's formula,
$$
    e^{tA}v_i = \sum_{k=1}^n \frac{t^AA^kv_i}{k!} = \sum_{k=1}^n \frac{t^A(a_i+b_i\-i)v_i}{k!} = e^{a_i+\-i b_i}v_i = e^{ta_i}(\cos tb_i + \-i\sin tb_i)v_i
    \xrightarrow{t\to \infty} 0\cdot v_i = 0,
$$
where the convergence follows from $a_i < 0$ and $|\cos \theta + \-i\sin\theta| = 1$ for any $\theta$.
Since $A$ is diagonalizable, $\Span(v_1,\ldots v_n) = \C^n$, which implies that $e^{tA} \to 0$ as $t\to \infty$.

Let $M$ be the maximum of real part of eigenvalues of $A$.
Since the diagonalizable matrices are dense in $\C^{n\times n}$, there is a sequence of diagonalizable matrices $(A_k')_{k\geq 1}$ that converges to $A$ with $\max_j\Re\lambda_j(A_k)\le M/2$ for any $k\in[n]$.
Hence, $\norm{e^{tA_k}}_2 \leq e^{tM/2}$ holds and the convergence is uniform.
We then can swap the limits by uniform convergence,
$$
    \lim_{t\to \infty} e^{tA} = \lim_{t\to \infty} \lim_{k\to \infty}e^{tA_k}
    = \lim_{k\to \infty}\lim_{t\to \infty} e^{tA_k} = 0 .
$$
\section*{Problem 3}
\textbf{(a)}
Since $v^\top A v = (v^\top A v)^\top = v^\top A^\top v$, we have that
$$
   v^\top A v = 0 \iff 2v^\top A v = v^\top A v + v^\top A^\top v = v^\top(A+A^\top)v = 0.
$$
Therefore, replacing $A$ with $A+A^\top$ does not change the problem.
\\\textbf{(b)}
Since $A^\top = A$, $A = U\Lambda U^\top = \sum_{i=1}^n \lambda_i u_iu_i^\top$ for orthonormal $U = (u_1,\ldots,u_n)$, where $\lambda_1\geq \ldots \geq \lambda_n$ and $\Lambda = \diag\{\lambda_i\}_{i\in[n]}$.
\\$\implies:$ Let $v$ be the non-trivial isotropic vector.
We have that 
$$
    v^\top Av = \sum_{i=1}^n \lambda_i \inner{v}{u_i}^2 = 0\implies (\forall i\in[n], \lambda_i\inner{v}{u_i}=0)\lor (\lambda_1 > 0\land \lambda_n < 0).
$$
Since $Av = \sum_{i=0}^n \lambda_i\inner{v}{u_i}u_i \neq 0$, we conclude that $\lambda_1 > 0$ and $\lambda_n < 0$.
\\$\impliedby:$ Define $u(t) = tu_1 + (1-t)u_n$.
For any $t\in[0,1]$, the following holds,
$$
    Au(t) = \lambda_1 tu_1 + \lambda_n(1-t)u_n \neq 0,
$$
since $u_1,\ldots,u_n$ are linearly independent and $\lambda_1 > 0,\lambda_n < 0$.
Let $f(t) = u(t)^\top A u(t)$.
Trivially $f$ is a quadratic function and $f$ is continuous.
Therefore, by $f(0) = \lambda_n < 0$, $f(1) = \lambda_1 > 0$ and continuity, there exists $t' \in (0,1)$ such that $f(t') = 0$.
That is, $u(t')^\top Au(t') = 0$ and $Au(t') \neq 0$.
Hence, $u(t')$ is the non-trivial isotropic vector.

\section*{Problem 4}
\section*{Problem 5}
\section*{Problem 6}

\end{document}